\documentclass[10pt]{article}
\usepackage[T1]{fontenc}
\usepackage{calc}
\usepackage{setspace}
\usepackage{multicol}
\usepackage{fancyheadings}

\usepackage{graphicx}
\usepackage{color}
\usepackage{rotating}
%\usepackage{aer}
%\usepackage{aertt}
\usepackage{verbatim}

\usepackage{hyperref}
\usepackage{xcolor}
\hypersetup{
    colorlinks,
    linkcolor={red!50!black},
    citecolor={blue!50!black},
    urlcolor={blue!80!black},
}
\newcommand{\bulurl}[1]{\url{#1}}


\usepackage{array}
\usepackage{ragged2e}

\setlength{\voffset}{0in}
\setlength{\topmargin}{0pt}
\setlength{\hoffset}{-0.25in}
\setlength{\oddsidemargin}{0pt}
\setlength{\headheight}{0pt}
\setlength{\headsep}{-0.3in}
\setlength{\marginparsep}{0pt}
\setlength{\marginparwidth}{0pt}
\setlength{\marginparpush}{0pt}
\setlength{\footskip}{0.3in}
\setlength{\textwidth}{7in}
\setlength{\textheight}{9.5in}
\setlength{\parskip}{0pc}
\setlength{\parindent}{0pc}

\renewcommand{\baselinestretch}{1.0}

\newcommand{\bi}{\begin{itemize}}
\newcommand{\ei}{\end{itemize}}
\newcommand{\be}{\begin{enumerate}}
\newcommand{\ee}{\end{enumerate}}
\newcommand{\bd}{\begin{description}}
\newcommand{\ed}{\end{description}}
\newcommand{\prbf}[1]{\textbf{#1}}
\newcommand{\prit}[1]{\textit{#1}}
\newcommand{\beq}{\begin{equation}}
\newcommand{\eeq}{\end{equation}}
\newcommand{\bdm}{\begin{displaymath}}
\newcommand{\edm}{\end{displaymath}}
\newcommand{\script}[1]{\begin{cal}#1\end{cal}}
\newcommand{\citee}[1]{\citename{#1} (\citeyear{#1})}
\newcommand{\h}[1]{\hat{#1}}
\newcommand{\ds}{\displaystyle}
\newcommand{\centerpic}[2]{\begin{center}\includegraphics[#1]{#2}\end{center}}

\newcommand{\T}{\rule{0pt}{3.1ex}}
\newcommand{\B}{\rule[-3.1ex]{0pt}{0pt}}
\newcommand{\TT}{\rule{0pt}{10ex}}
\newcommand{\BB}{\rule[-15ex]{0pt}{0pt}}

\newcommand{\toprule}{\par\vspace*{5pt}\noindent{\hrule\hfill}\par\vspace*{1pt}}
\newcommand{\botrule}{\par\noindent{\hrule\hfill}\par}

\pagestyle{plain}

\begin{document}

\begin{center} 
\textbf{Course \#\# and Title - Semester Year}\\
\end{center} 

\botrule \textbf{Instructor Information} \toprule
James Murray, Ph.D.\newline
Office: 403T Wimberly Hall\newline
Email: jmurray@uwlax.edu\newline
Office Phone: 608-785-5140 \newline

\botrule \textbf{Meeting Time / Location} \toprule
\ \\

\botrule \textbf{Course Goals} \toprule
\ \\

\botrule \textbf{Learning Objectives} \toprule
\ \\

\botrule \textbf{Course Resources} \toprule
\bd
\item Textbook: 
\item Class websites: \bulurl{http://www.murraylax.org/}.  All material handed out in class will be posted on the class website.  
\item Desire2Learn (D2L):  Grades will be posted on D2L, and some quizzes may be administered on D2L.  Most of the class material will be posted on the class website. \\ 
\ed

\botrule \textbf{Office Hours} \toprule
I am available for office hours \textbf{by appointment with a minimum of only one hour notice.}  You may schedule a 15 minute appointment by visiting \bulurl{https://murraylax.youcanbook.me}.  The blocks of time that I am available each week vary and are kept up to the minute on the YouCanBook.Me online scheduler.  Additional walk-in office hours will be added as necessary, especially during exam weeks and weeks with significant homework deadlines.  My typical weekly availability is given below. \\

\hspace*{-0.5pc}\begin{tabular}{lcl}
\hspace*{-0.5pc}\begin{tabular}{ll}
  10:00 AM - 11:30 AM & Monday through Thursday \\
\end{tabular}

& &

\begin{tabular}{p{2.5in}cl}
Scan code or visit\newline \bulurl{https://murraylax.youcanbook.me} \newline to make an office hours appointment.
& ~~ &
\raisebox{-0.75\height}{\includegraphics[scale=0.25]{officehours.png}}
\end{tabular}

\end{tabular}\\

Office hours are not a substitute for attending class.  Except when missing class for very extreme circumstances that were promptly discussed with me, it is not acceptable to use office hours to ask questions about material you missed while not in attendance.\\

\botrule \textbf{OFFICE HOURS SATISFACTION GUARANTEED!} \toprule
I guarantee that coming to office hours to ask questions on material you do not completely understand will increase your grade in the class.  If you come to office hours and still get questions wrong on the exam about the topic we discussed, I will refund those points!  This guarantee is subject to the following conditions:
\be
\item Refund points are good for up to two topics on each exam, and up to a total of 10 percentage points per exam.
\item You and I must agree in writing using the form at the end on the syllabus that the topic that is guaranteed.  On this form we both agree that I have explained the topic sufficiently and that you have acquired a good understanding from our discussion.  
\item You must have attended class when the topic was initially taught.
\item Only three guarantee forms may filed before each exam.
\item Only one guarantee form may be filed on a single day.  Therefore, filing two or three guarantee forms before an exam will require at least two or three separate office hours visits, respectively.
\item After the exam, if you get questions wrong that were guaranteed, it is your responsibility to recognize this and present your copy of the guarantee form to notify me you should get the points refunded.
\ee

\botrule \textbf{Assessment} \toprule
\ \\

\botrule \textbf{Grade Breakdown} \toprule
\noindent\begin{tabular}{cccc} \\
94-100 & A & 77-81 & BC  \\
89-93 & AB & 70-76 & C \\
82-88 & B & 0-69 & F \\
\end{tabular}\\

I reserve the right to scale every person's grade up by the same amount on any graded item in the event that much of the class falls short of the scale above.  However, even if the grades are significantly low, there is no guarantee that I will ever do this.  \\

\botrule \textbf{In-class Exercises and Quizzes} \toprule
\ \\  

\botrule \textbf{Homework Assignments} \toprule
\ \\

\botrule \textbf{Writing Assignments} \toprule
\ \\

\botrule \textbf{Projects} \toprule
\ \\

\botrule \textbf{Exams} \toprule
There will be two midterm exams and one final exam.  The midterm exams cover primarily material since the previous exam, but the second exam may include earlier material.  The final exam is cumulative.  Unless otherwise enforced by university policy, you must take the final exam at the scheduled given below.  The exam dates are as follows: 
\bi
\item Exam 1: Day, Date
\item Exam 2: Day, Date
\item Final Exam: 
The UWL final exam schedule can be found here:\\ \bulurl{http://www.uwlax.edu/Records/Final-Exam-Schedule/}
\ei

Exams may require a calculator, so be sure to bring a calculator to class on exam days.  You may use a basic or scientific calculator, but you are not permitted to use electronic devices with capacity to store information.  Therefore, devices such as graphing calculators, mobile phones, tablets, and electronic translators are not permitted.  If English is not your first language, you are permitted to bring a paperback English translator/dictionary.  You may not share calculators with other students. \\

\botrule \textbf{Graded Coursework Response Time} \toprule
I will return all graded work to you within two weeks of the due date, or before the next exam date if the work is relevant practice for the exam.  Work that is turned in late may not be accepted, but if so, I may not be able to grade the work within the two-week time frame.  Grades will be posted on the D2L gradebook on or before the date I return the graded work to the class.  I will return graded coursework in compliance with FERPA regulations, such as in class or during my office hours.  I will bring your graded coursework to class only once.  If you are not in attendance when coursework is returned, it is your responsibility to make arrangements to pick your work.\\

\botrule \textbf{Grade Change or Extra Credit Requests} \toprule 
I evaluate and assign grades for a lot of work from a lot of students in multiple courses, so it possible or even likely I may make mistakes.  It is appropriate and helpful to me if you keep track of your grades that I post to D2L and notify me if I have made a mistake.  \textbf{It is not appropriate to ask for grade changes or special extra credit opportunities after performing poorly on assignments or exams}, or not achieving a grade that you hoped to earn.  These requests will not be granted.  Know also that such requests to me and most any instructor reflect poorly on your professionalism, attitude, and priorities.\\

In the event that scores are low for an assignment or exam across most or all students I do reserve the right to increase every students' grade by the same amount or give an extra credit opportunity to all students.  Please understand that such events are rare and will probably not happen in a typical semester.  Please do not make requests for such grade changes or extra credit opportunities.  I do see the full distribution of all the students' grades and I have over a decade of knowledge and experience of the historical performances of students on similar assignments and exams in my classes and with other instructors.  I use this full range of information to determine when such changes are appropriate.\\

\botrule \textbf{Double Counting Work with Other Courses} \toprule
You may be allowed to use an assignment from this course to also satisfy a requirement for another course if you adhere to the following requirements:
\be
\item Notify both me and the instructor of your other course of your intention in writing in a single email that is addressed to us both.
\item Explain how your idea satisfies the assignment requirements for both courses. You should review the assignment guidelines and rubrics from both courses, if applicable.
\item Clearly articulate how the one assignment will result in something more valuable and more sophisticated than if you completed the two assignments separately.  You need to establish here that by completing a single assignment to satisfy both courses' requirements results in \textit{higher quality work and not less total work}.
\item Both instructors give written permission (an email will suffice) which includes a notification to the other instructor.
\ee

Submitting the same work in multiple courses without obtaining permission from both instructors is considered self-plagiarism.  Self-plagiarism is a type of plagiarism; it is considered academic dishonesty and it will be treated as such.\\

\botrule \textbf{Submitting Late Work} \toprule
Attendance is required to receive credit for graded work completed in class, including quizzes, in-class exercises, and exams.  Late work will not be accepted except for extraordinary circumstances.  If possible, you should notify me before a missed class or assignment deadline that you will not be able to complete the assigned work at the given deadline and we may agree on an appropriate accommodation.  In the event of missing class or an assignment deadline due to illness or emergency, you should notify me as soon as possible afterwards so that we can make arrangements for you to make up missed work.  \textbf{Late work cannot be accepted under any circumstances after I have returned graded work to the class or after answer keys have been circulated.}  It is for this reason that is important to notify me as soon as possible any circumstance that makes it appropriate to submit work late.\\

\botrule \textbf{Attendance} \toprule
Full physical and mental attendance is expected for every class period. That means \textbf{you attend and are attentive} in every class, and you take responsibility for knowing any announcement or concept I describe in class.  If you need to miss a class day or exam day because of illness or emergency you may be excused and arrangements can be made for you to make up missed work if you notify me as soon as possible after the missed class.  \\

\botrule \textbf{Eagle Alert System} \toprule 
This class will be participating in the UWL Eagle Alert system through WINGS. The Early Alert system is designed to promote student success. If I notice that you are experiencing difficulties early in the semester (e.g., low assignment scores, poor attendance, minimal engagement in the classroom), I may enter feedback into the program and you will receive an email indicating that feedback has been left.  I may also enter positive feedback encouraging you to think about additional opportunities. You will be able to access the feedback through your student center in WINGS.  I encourage you to meet with me and use one or more of several helpful campus resources listed here \bulurl{http://www.uwlax.edu/studentsuccess/}. \\

\botrule \textbf{Online Student Evaluation of Instruction (SEI)} \toprule
The university conducts student evaluations electronically. Approximately 2 weeks prior to the conclusion of a course, you will receive an email at your UWL email address directing you to complete an evaluation for each of your courses. In-class time will be provided for students to complete the evaluation in class. Electronic reminders will be sent if you do not complete the evaluation. The evaluation will include numerical ratings and, depending on the department, may provide options for comments. The university takes student feedback very seriously and the information gathered from student evaluations is more valuable when a larger percentage of students complete the evaluation. Please be especially mindful to complete the surveys.\\

\botrule \textbf{Sexual Misconduct} \toprule
As an employee of the University of Wisconsin-La Crosse, I am a mandated reporter of sexual harassment and sexual violence that takes place on campus or otherwise affects the campus community. This means that if I receive detailed or specific information about an incident such as the date, time, location, or identity of the people involved, I am obligated to share this with UWL's \href{https://www.uwlax.edu/affirmative-action/}{Title IX Coordinator} (\bulurl{https://www.uwlax.edu/affirmative-action/}) in order to enable the university to take appropriate action to ensure the safety and rights of all involved. For students not wishing to make an official report, there are confidential resources available to provide support and discuss the available options. The contact in Student Life is Ingrid Peterson, Violence Prevention Specialist, 608.785.8062, \texttt{\href{mailto:ipeterson@uwlax.edu}{ipeterson@uwlax.edu}}. Please see \bulurl{http://www.uwlax.edu/sexual-misconduct} for more resources or to file a report.\\

\botrule \textbf{Religious Accommodations} \toprule
Per the UWL Undergraduate and Graduate \href{http://catalog.uwlax.edu/undergraduate/aboutuwlax/}{Catalogs} (\bulurl{http://catalog.uwlax.edu/undergraduate/aboutuwlax/}), ``any student with a conflict between an academic requirement and any religious observance must be given an alternative means of meeting the academic requirement. The student must notify the instructor within the first three weeks of class (within the first week of summer session and short courses) of specific days/dates for which the student will request an accommodation. Instructors may schedule a make-up examination or other academic requirement before or after the regularly scheduled examination or other academic requirement.''\\

\botrule \textbf{Students with Disabilities} \toprule
Any student with a documented disability (e.g. ADHD, Autism Spectrum Disorder, Acquired Brain Injury, PTSD, Physical, Sensory, Psychological, or Learning Disability) who needs to arrange academic accommodations should contact The ACCESS Center (165 Murphy Library, 608-785-6900, \texttt{\href{mailto:ACCESSCenter@uwlax.edu}{ACCESSCenter@uwlax.edu}}) and meet with an adviser to register and develop an accommodation plan.  In addition to registering with The ACCESS Center, it is the student's responsibility to discuss their academic needs with their instructors.  You can find out more about services available to students with disabilities at \href{http://www.uwlax.edu/access-center}{The ACCESS Center website}: \bulurl{http://www.uwlax.edu/access-center}.\\

\botrule \textbf{Veterans and Active Military Personnel} \toprule
Veterans and active military personnel with special circumstances (e.g., upcoming deployments, drill requirements, disabilities) are welcome and encouraged to communicate these, in advance if possible, to me. For additional information and assistance, contact the \href{http://www.uwlax.edu/veteran-services/}{Veterans Services Office} (\bulurl{http://www.uwlax.edu/veteran-services/}). Students who need to withdraw from class or from the university due to military orders should be aware of the \href{http://catalog.uwlax.edu/undergraduate/academicpolicies/withdrawal/\#military-duty-withdrawal-university}{military duty withdrawal policy} (\bulurl{http://catalog.uwlax.edu/undergraduate/academicpolicies/withdrawal/\#military-duty-withdrawal-university}).\\

\botrule \textbf{E-mail Guidelines} \toprule
I insist on the following e-mail etiquette rules (many of these are also recommended by the College of Business Administration).  Failure to adhere to these guidelines will result in a reply with a friendly reminder to follow these e-mail guidelines.
\bi
\item Allow one business day to elapse before expecting a reply. 
\item Questions sent by e-mail should be able to be answered with only a few words, such as 'yes' or 'no' questions.  Questions whose answers involve explaining class material are not appropriate over e-mail.  For answers to these questions you should make an office hours appointment.
\item Always include a subject that is brief but still has sufficient detail that includes this course number.
\item Look at your class notes and syllabus before sending an e-mail.  Do not ask a question whose answer is on the syllabus or announced in class (unless you missed class for a legitimate reason).
\item Always spell check, grammar check, and re-read your e-mail before sending it.
  \ei

\botrule \textbf{Academic Misconduct} \toprule
Academic misconduct is a violation of the UWL Student Honor Code and is unacceptable. I expect you to submit your own original work and participate in the course with integrity and high standards of academic honesty.  The UWL Student Honor Code can be found online at \bulurl{http://catalog.uwlax.edu/undergraduate/academicpolicies/studentconduct/}.  In the event a student is caught committing academic misconduct, I will pursue the harshest penalties allowed according to the UWS 14 Student Academic Disciplinary Procedures, which can be found here \bulurl{http://www.uwlax.edu/Student-Life/Student-handbook/\#14}.\\

\botrule \textbf{UWL Legal Obligations to Students} \toprule
Many of the above legal obligations of instructors and the university are afforded to all UWL students in all their courses.  You can see these at \bulurl{https://www.uwlax.edu/info/syllabus/}.\\

\botrule \textbf{Topics Schedule} \toprule
 Below is a list of topics and textbook readings for this class.    \toprule

\newpage

\thispagestyle{empty}

\begin{center}[MORE COPIES ARE AVAILABLE OUTSIDE OF 403T WIMBERLY]\end{center}

\toprule
\begin{center}\Large{\textbf{Office Hours Guarantee}}\end{center}
\vspace*{0.1in}

\begin{tabular}{ll}
\textbf{Name:}  & \line(1,0){200} \\ [1pc]
\textbf{Date / Time:} & \line(1,0){200} \\  [1pc]
\textbf{General Topic:}  & \line(1,0){200} \\  [1pc]
\textbf{Chapter / Reading Reference:}  & \line(1,0){200}\\  [1pc]
\end{tabular}

\begin{tabular}{l}
\textbf{Describe specifically the concept that is guaranteed:}
\end{tabular}
\vspace*{1.6in}

\begin{tabular}{cc}
\line(1,0){200} & \line(1,0){200} \\
Student Signature & Instructor Signature \\
\end{tabular} \\\\

Student Copy 
\vspace*{0.1in}

\toprule
\begin{center}\Large{\textbf{Office Hours Guarantee}}\end{center}
\vspace*{0.1in}

\begin{tabular}{ll}
\textbf{Name:}  & \line(1,0){200} \\ [1pc]
\textbf{Date / Time:} & \line(1,0){200} \\  [1pc]
\textbf{General Topic:}  & \line(1,0){200} \\  [1pc]
\textbf{Chapter / Reading Reference:}  & \line(1,0){200}\\  [1pc]
\end{tabular}

\begin{tabular}{l}
\textbf{Describe specifically the concept that is guaranteed:}
\end{tabular}
\vspace*{1.6in}

\begin{tabular}{cc}
\line(1,0){200} & \line(1,0){200} \\
Student Signature & Instructor Signature \\
\end{tabular} \\\\

Instructor Copy 
\vspace*{0.1in}
\toprule
\end{document}



